\section{Discussion and Limits}
\label{sec:discussion}
\subsection{When the tradeoff is useful}
The strongest case for cooperative analysis is a caller that values reducing
concentrated engine work and can tolerate an additional hop of spectrum age.
The representative comparison reduces callback p99 while increasing aggregate
cost. That is useful evidence for an engineering choice, not a new transform
complexity result. For a low-cost visualizer, headroom in the callback budget
may matter more than minimum arithmetic throughput; an application needing
fresh output or many concurrent analyzers may make the opposite choice.
The observed timing distributions cannot decide those application priorities.
For the representative $D=64$, 48\,kHz case, the nominal callback interval is
\PaperCmpBudgetUs\,$\mu$s: the production and vDSP p99 values occupy only
\PaperCmpCoreBudgetPercent\% and \PaperCmpVdspBudgetPercent\% of that interval,
respectively. Thus the measured reduction does not show that the native
baseline was near its deadline, or that redistribution prevented audible
failures. Such a claim would require a contended host and device measurements.

The PFFFT hybrid is a useful intermediate design: it distributes surrounding
work while keeping a native FFT indivisible. It shows that scheduling the
pipeline around a batch transform can capture some benefit without suspending
individual butterflies. Its comparison with the matched scheduled-batch
control isolates work placement more closely than a comparison with the
production analyzer, which also changes arithmetic, layout, storage, and
caching. The extra state in a resumable transform is justified only when its
measured behavior meets a requirement that the simpler alternative does not.

The inverse and chain results sharpen that qualification. Deferring butterfly
execution does not eliminate the bulk buffering and conjugation at job or
stage boundaries. The representative native PFFFT and vDSP inverse jobs provide both lower
aggregate cost and lower p99 than the incremental control. The complete-chain
control reduces p99 more modestly than its large aggregate-cost penalty. Thus the favorable
analysis example cannot be generalized to all frequency-domain processing.
The formulas in Section~\ref{sec:inverse-chain} identify what a future complete
inverse scheduler would need to change. The current measurements are already
useful as counterexamples to a blanket claim that incremental execution
reduces callbacks more effectively than optimized batch execution.

\subsection{Interpretation and external validity}
The primary tables compare declared input/output contracts. Smoothing
interval arithmetic and transform layouts differ across implementation
families; numerical agreement is within the stated contracts and budgets,
not bitwise identity. Precision, channel count, state, callback phase, and
additional DSP change the task. Extension results should be read within each
matched group rather than combined into a single provider score.

A high percentile, a maximum, and aggregate throughput answer different
questions. Single-sample p99 can omit frame work occurring less than once per
hundred calls. Maxima retain those bursts and rare preemptions but depend on
observation duration; the separate group durations preclude indiscriminate
maxima comparisons. No samples are trimmed and no empty-timer cost is
subtracted. Small absolute differences and sub-tick observations should not
be interpreted as finely resolved instruction costs. Storage probes measure
only their declared allocation scope; unknown provider-native storage prevents
a complete memory ranking.

All confirmation observations come from one compiler and one macOS ARM64 host
on one night. The prepared sessions expose temporal variation, including
unfavorable session excursions, but do not establish independent-host
replication or a stable rare-event distribution. Dynamic FFTW planning is part
of the recorded configuration. Other architectures, compilers, and plans may
change the relative costs. Binary64 references, deterministic fixtures, and
all-publication audits support the stated numerical workloads, not every
possible input or application-specific tolerance.

\subsection{Work bounds and real-time behavior}
The proof bounds credits in the complete analysis sequence. Unequal unit
costs, memory hierarchy, preemption, host lifecycle work, and unrelated DSP
still determine elapsed time. A late call executes its logical quota; the
scheduler does not recover a missed deadline. Producer publication also does
not guarantee display consumption or repaint cadence: a slow display may skip
complete snapshots by design.

The benchmark runs synchronously without an audio driver or concurrent UI.
It therefore provides no measured device underrun rate or proof of hard
real-time safety. Worker-thread analyzers, overlap-reusing sliding/hopping
methods, non-power-of-two transforms, and long-filter applications are not
ranked. Appendix~\ref{app:evaluation-agenda} identifies the next experiments.
The complete paper can support an implementation decision under its recorded
conditions; broader recommendations require those additional measurements.
